% Generated draft: requires mathematical and source review.
\documentclass[11pt]{article}
\usepackage[margin=25mm]{geometry}
\usepackage{fontspec}
\setmainfont{Latin Modern Roman}
\usepackage{amsmath,amssymb,mathtools,unicode-math}
\setmathfont{Latin Modern Math}
\usepackage{enumitem,needspace,xurl,hyperref}
\setlength{\parindent}{0pt}
\setlength{\parskip}{5pt}
\setlength{\emergencystretch}{3em}
\newcommand{\N}{\mathbb{N}}
\newcommand{\Z}{\mathbb{Z}}
\newcommand{\Q}{\mathbb{Q}}
\newcommand{\R}{\mathbb{R}}
\newcommand{\C}{\mathbb{C}}
\newcommand{\F}{\mathbb{F}}
\newcommand{\A}{\mathbb{A}}
\newcommand{\PP}{\mathbb{P}}
\newcommand{\E}{\mathbb{E}}
\newcommand{\eps}{\varepsilon}
\newcommand{\dd}{\,\mathrm{d}}
\begin{document}
\section*{UnsolvedMath ID: 30004750}
\textbf{Source problem code:} catalogue entry
\subsection{Definitions and mathematical statement}
A number field is a finite extension \(K/\Q\). Its ordinary ideal class group is
\[
\operatorname{Cl}(K)=I_K/P_K,
\]
where \(I_K\) is the group of nonzero fractional ideals of the ring of integers \(\mathcal O_K\), and \(P_K\) is its subgroup of nonzero principal fractional ideals. This recalls the ordinary class-group convention; the supplied record does not specify whether the intended heuristics concern ordinary class groups, another class-group variant, or particular primary components.

The notation \(n\)-Selmer refers to Selmer theory associated with multiplication by a positive integer \(n\), in a setting where that construction is defined. A precise definition requires hypotheses on the algebraic group and its base, together with a cohomological construction and local conditions. The source does not supply these choices for its proposed class of possibly nonproper algebraic groups. Consequently, no single definition of an \(n\)-Selmer group for the entire intended scope is imposed here. Likewise, distributional heuristics require specified families, an ordering or probability measure, and limiting predictions; these are not supplied.

\textbf{Research question.} Can one formulate unified heuristics for \(n\)-Selmer groups of possibly nonproper algebraic groups that also imply the Cohen--Lenstra--Martinet heuristics for class groups of number fields?

The requested output is a coherent formulation of heuristics and an explanation of how the relevant Cohen--Lenstra--Martinet predictions follow from it. The source does not specify the allowed indices \(n\), the base fields, the class of algebraic groups, the local conditions, the varying families, or the precise versions of the Cohen--Lenstra--Martinet predictions to be recovered. Thus this is a research-program question whose full quantifiers and hypotheses remain to be formulated, rather than a completely specified conjectural assertion.
\subsection{Short English statement}
Seek a common framework predicting the distributions of \(n\)-Selmer groups for algebraic groups, including nonproper groups, that also yields the Cohen--Lenstra--Martinet class-group heuristics. The source leaves the groups, fields, local conditions, families, and precise predictions unspecified.
\subsection{Sources}
\begin{itemize}
\item Local UnsolvedMath registry: permanent ID \(30004750\), source problem code \texttt{catalogue entry}, title \emph{Unified Selmer and Cohen-Lenstra-Martinet Heuristics}; collection: Oberwolfach Reports Open Problems. No individual proposer is identified. The supplied proposed year is 2021.
\item Supplied report citation: \emph{Explicit Methods in Number Theory (2022)}, Oberwolfach Reports, pp.\ 1805--1858, DOI \texttt{10.4171/owr/2021/34}. \url{https://doi.org/10.4171/owr/2021/34}. The supplied background describes Oberwolfach Reports as documenting workshops and their open-problem discussions.
\item Supplied literature source, labelled ``Landesman'': \url{https://arxiv.org/abs/2106.10357}. No further bibliographic details or precise supporting passage are supplied.
\item Supplied literature source, labelled ``Oberwolfach report'': \url{https://ems.press/content/serial-article-files/46913}.
\item Literature assessment dated exactly \texttt{2026-08-22}: the supplied curation labels the problem open and states that unified heuristics for Selmer groups of general algebraic groups, with the requested implication for Cohen--Lenstra--Martinet heuristics, remain a research question. This is supplied dated curation, not an independently verified literature assessment; the source explicitly calls for independent verification of time-sensitive status and cited claims.
\item The registry labels the statement \texttt{corrected\_verified} and reports manual semantic, typographic, and LaTeX review. These are supplied editorial claims, not independent verification here of mathematical precision or current status.
\end{itemize}
\end{document}
